\section{Exact trial vectors for opposite fields}\label{app:trials}

We prove the four variational inequalities used to initialize the
boundary argument.  Throughout this appendix,
\[
 a=\frac{700741}{500000},\qquad h=\frac35,\qquad
 \alpha=\frac{2021}{5000}=.4042,
\]
and the unshifted opposite-field Hamiltonian on $m$ sites is
\begin{equation}\label{apptr:hamiltonian}
 H_m^{\mathrm{opp}}
 =\sum_{j=1}^{m-1}\mathbf S_j\cdot\mathbf S_{j+1}
   -\frac35S_1^z+\frac35S_m^z.
\end{equation}
The bulk matrices below are those of the periodic variational
construction in~\cite[Section ``Two variational inputs'']{periodic}.
We give their complete integer specification and the new endpoint
vectors, so the present finite estimates can be reproduced without
importing a variational bound for open chains.

\begin{proposition}\label{apptr:bound}
Let $E_0^{\mathrm{opp}}(m)$ be the smallest eigenvalue of
$H_m^{\mathrm{opp}}$.  Then
\begin{equation}\label{apptr:four-bounds}
 E_0^{\mathrm{opp}}(m)+am<\alpha
 \qquad\text{for }m\in\{34,78,142,240\}.
\end{equation}
\end{proposition}

The proof consists of a fixed integer matrix-product vector, an exact
contraction for its norm and energy, and four integer comparisons.
No spectral approximation enters the calculation.

\subsection{The complete tensor and endpoint data}

Let $\mathcal V$ have the $26$ labels
\begin{equation}\label{apptr:virtual-labels}
 (i,d),\qquad 1\le i\le8,\qquad
 d=-\ell_i,-\ell_i+2,\ldots,\ell_i,
 \qquad(\ell_i)_{i=1}^8=(1,1,1,1,3,3,3,5).
\end{equation}
Order them first by $i$ and then by increasing $d$.  Three real
$26$-by-$26$ matrices $A_s$, $s\in\{-1,0,1\}$, are specified as follows.
At row $(i,d)$ and column $(j,e)$ put
$\ell=\ell_i$, $k=\ell_j$, and $r=(\ell+d)/2$.
The entry is zero unless
\begin{equation}\label{apptr:support}
 e=d+2s,\qquad |k-\ell|\le2,\qquad
 k=\ell\ \Longrightarrow\ i=j.
\end{equation}
For an allowed entry, multiply the coefficient in the second column of
Table~\ref{apptr:entry-table} by the factor in its $s$ column.
The parameters $D_i$ and $P_{ij}$ appear in
Tables~\ref{apptr:endpoint-table} and~\ref{apptr:coupling-table}.
The latter parameter is defined only when $\ell_j=\ell_i+2$.
All labels in these rules must belong to~\eqref{apptr:virtual-labels};
an out-of-range column is absent.

Define columns $\mathbf b,\mathbf c\in\mathbb R^{26}$ by
\begin{equation}\label{apptr:endpoints}
 \mathbf b_{(i,d)}=\begin{cases}\beta_i,&d=-1,\\0,&d\ne-1,\end{cases}
 \qquad
 \mathbf c_{(i,d)}=\begin{cases}\kappa_i\beta_i,&d=-1,\\0,&d\ne-1.\end{cases}
\end{equation}
All their parameters are included in Table~\ref{apptr:endpoint-table}.

For each $m\ge2$, the physical vector is
\begin{equation}\label{apptr:physical-vector}
 \psi_m=\sum_{s_1,\ldots,s_m=-1}^{1}
 \bigl(\mathbf b^{\mathsf t}A_{s_1}\cdots A_{s_m}\mathbf c\bigr)
 e_{s_1}\otimes\cdots\otimes e_{s_m}.
\end{equation}
This is the open-boundary matrix-product form of
\cite[Section 3.1, Equation (11)]{PerezGarcia2007}, with endpoint rows
and columns absorbed into the first and last matrices.
The vectors $e_{-1},e_0,e_1$ are the orthonormal spin basis fixed in
the statement of the problem.  In particular, the coefficients in
\eqref{apptr:physical-vector} are ordinary real amplitudes: no virtual
inner product or physical basis weights are inserted.  Normalization
of the matrices $A_s$ is unnecessary.  The support rule forces a
nonzero amplitude to have $\sum_j s_j=0$, because both endpoints have
virtual weight $-1$.  This restriction on the trial vector does not
restrict the variational upper bound on the smallest eigenvalue of
the full Hamiltonian.

The following tables complete the integer specification of these matrices
and endpoints.

\begin{table}[htbp]
\centering
\setlength{\tabcolsep}{5pt}
\begin{tabular}{c|c|ccc}
 &coefficient&$s=-1$&$s=0$&$s=1$\\ \hline
 $k=\ell+2$&$P_{ij}$&
 $2(\ell+1-r)(\ell+2-r)$&$2(r+1)(\ell-r+1)$&$(r+1)(r+2)$\\
 $k=\ell$&$D_i$&$2(\ell-r+1)$&$d$&$-(r+1)$\\
 $k=\ell-2$&$-P_{ji}$&$2$&$-2$&$1$
\end{tabular}
\caption{Every allowed entry of $A_s$ is the coefficient times the
listed factor.  All other entries are zero.}
\label{apptr:entry-table}
\par\vspace{1.5em}
\begin{tabular}{r|r|r|r|r}
 $i$&$\ell_i$&$D_i$&$\beta_i$&$\kappa_i$\\ \hline
 1&1&$-289$&10000&12\\
 2&1&$-84$&5620&12\\
 3&1&$-9$&2456&12\\
 4&1&33&4868&12\\
 5&3&$-22$&82&6\\
 6&3&27&$-809$&6\\
 7&3&71&12475&6\\
 8&5&$-11$&$-9645$&1
\end{tabular}
\caption{Virtual types, diagonal tensor parameters, and endpoint
parameters.  The left and right endpoints are supported on $d=-1$.}
\label{apptr:endpoint-table}
\par\vspace{1.5em}
\begin{tabular}{c|r@{\qquad}c|r@{\qquad}c|r}
 $(i,j)$&$P_{ij}$&$(i,j)$&$P_{ij}$&$(i,j)$&$P_{ij}$\\ \hline
 $(1,5)$&$-11$&$(2,7)$&$-49$&$(4,6)$&9\\
 $(1,6)$&$-35$&$(3,5)$&$-6$&$(4,7)$&$-32$\\
 $(1,7)$&$-28$&$(3,6)$&$-25$&$(5,8)$&1\\
 $(2,5)$&$-23$&$(3,7)$&$-20$&$(6,8)$&2\\
 $(2,6)$&16&$(4,5)$&$-34$&$(7,8)$&$-16$
\end{tabular}
\caption{All fifteen couplings with $\ell_j=\ell_i+2$.}
\label{apptr:coupling-table}
\end{table}
\FloatBarrier

\subsection{Norm, fields, and the bond insertion}

For $s,t\in\{-1,0,1\}$, put $B_{st}=A_sA_t$ and define
\begin{align}
 T_*&=\sum_{s=-1}^{1}A_s\otimes A_s,
 &G_*&=\sum_{s=-1}^{1}sA_s\otimes A_s,
 \label{apptr:one-site-transfers}\\
 O_*&=\sum_{s,t=-1}^{1}st\,B_{st}\otimes B_{st}
 \notag\\[-2pt]
 &\quad+
 \sum_{\substack{s=-1,0\\t=0,1}}
 \bigl(B_{st}\otimes B_{s+1,t-1}
       +B_{s+1,t-1}\otimes B_{st}\bigr).
 \label{apptr:bond-insertion}
\end{align}
We use the usual Kronecker convention
$(C\otimes D)_{(v,w),(v',w')}=C_{v,v'}D_{w,w'}$.
Write
\begin{equation}\label{apptr:boundary-tensors}
 L_0=(\mathbf b\otimes\mathbf b)^{\mathsf t},
 \qquad u=\mathbf c\otimes\mathbf c,
 \qquad L_i=L_0T_*^i.
\end{equation}

\begin{lemma}\label{apptr:contractions}
With the data above, the norm $V_m=\langle\psi_m,\psi_m\rangle$
and the scaled shifted energy numerator
\[
 H_m^*=500000\,
 \langle\psi_m,(H_m^{\mathrm{opp}}+am\Id)\psi_m\rangle
\]
satisfy
\begin{align}
 V_m&=L_mu,\label{apptr:norm}\\
 H_m^*&=500000\sum_{r=0}^{m-2}
          L_0T_*^rO_*T_*^{m-2-r}u
       -300000L_0G_*T_*^{m-1}u\notag\\
 &\quad+300000L_{m-1}G_*u+700741mL_mu.
 \label{apptr:energy-expanded}
\end{align}
In particular, $V_m$ and $H_m^*$ are integers.
\end{lemma}

\begin{proof}
Expanding $T_*^m$ sums the products of two identical real amplitudes
over all physical words.  This proves~\eqref{apptr:norm} directly in
the orthonormal physical basis.  Inserting $G_*$ at one site instead
of $T_*$ multiplies the summand by the weight $s$ and therefore
inserts $S^z$ at that site.

For a bond, the spin normalization gives
\[
 \mathbf S\cdot\mathbf S
 =S^z\otimes S^z+\tfrac12
   (S^+\otimes S^-+S^-\otimes S^+).
\]
Its diagonal entry on $(s,t)$ is $st$.  Its off-diagonal entries are
$1$ between $(s,t)$ and $(s+1,t-1)$ for
$s\in\{-1,0\}$ and $t\in\{0,1\}$, since each allowed ladder
entry is $\sqrt2$.  The two summands in each parenthesis
of~\eqref{apptr:bond-insertion} give the two directions of that
matrix element.  Thus $O_*$ inserts exactly one unit-coupling bond,
with no extra factor or metric.  Its position after $r$ sites gives
the $r$th summand of~\eqref{apptr:energy-expanded}.

The left and right endpoint insertions have coefficients
$500000(-3/5)=-300000$ and $500000(3/5)=300000$, respectively.
The constant shift contributes $500000am=700741m$ times the norm.
There are $m-1$ bond positions and no closing bond, giving the stated
formula.  Every matrix, row, and column in that formula is integral.
\end{proof}

\subsection{The invariant 162-dimensional contraction}

All contractions in Lemma~\ref{apptr:contractions} can be carried out
in a single $162$-dimensional space.  On a paired virtual label
$((i,d),(j,e))$, consider the difference $d-e$.  Each nonzero entry
of $A_s$ changes virtual weight by $2s$ from row to column.
Consequently every summand $A_s\otimes A_s$ of $T_*$ or $G_*$
preserves this difference.  A nonzero entry of $B_{st}$ changes
weight by $2(s+t)$.  The two factors in every tensor summand of
$O_*$ have the same total physical weight, including the pair
$(s,t)$ and $(s+1,t-1)$ in the off-diagonal terms.  Each such
summand also preserves $d-e$.

Let $\mathcal K_0$ be the coordinate subspace with paired labels
\begin{equation}\label{apptr:restricted-labels}
 \mathcal Q=\{((i,d),(j,d)):(i,d),(j,d)\in\mathcal V\}.
\end{equation}
The multiplicities of the virtual weights
$d=-5,-3,-1,1,3,5$ are $1,4,8,8,4,1$, respectively; hence
\begin{equation}\label{apptr:restricted-dimension}
 \dim\mathcal K_0=1^2+4^2+8^2+8^2+4^2+1^2=162.
\end{equation}
The preceding argument shows that $T_*,G_*,O_*$ are block diagonal
by weight difference.  Both $L_0$ and $u$ are supported in
$\mathcal K_0$, because the two endpoint vectors are supported on
$d=-1$.  We may therefore restrict these three matrices to
\eqref{apptr:restricted-labels} before multiplying them, taking
powers, or forming the expressions of
Lemma~\ref{apptr:contractions}.  In particular, squaring the
restricted $T_*$ gives the restriction of $T_*^2$; there is no
contribution through a discarded state.  From now on the same
symbols denote these restricted integer matrices and endpoint
tensors, ordered by the inherited lexicographic order.

\subsection{An exact row recurrence}

The energy sum in~\eqref{apptr:energy-expanded} need not be evaluated
separately at every bond.  Define integer rows by
\begin{align}
 L_i&=L_{i-1}T_*\quad(i\ge1),\notag\\
 N_1&=-300000L_0G_*,\notag\\
 N_i&=N_{i-1}T_*+
       L_{i-2}\bigl(500000O_*+700741T_*^2\bigr)
       \quad(i\ge2).
 \label{apptr:recurrence}
\end{align}
Then
\begin{equation}\label{apptr:energy-recurrence}
 V_m=L_mu,\qquad
 H_m^*=\bigl(N_m+300000L_{m-1}G_*+700741L_m\bigr)u.
\end{equation}
To verify the latter identity, induction on $m$ in
\eqref{apptr:recurrence} gives
\begin{equation}\label{apptr:N-expansion}
 N_m=-300000L_0G_*T_*^{m-1}
   +\sum_{r=0}^{m-2}
      L_0T_*^r(500000O_*+700741T_*^2)T_*^{m-2-r}.
\end{equation}
Its bond contributions and left field are exactly those in
\eqref{apptr:energy-expanded}.  Each of its $m-1$ shift
contributions equals $700741L_m$.  The final term $700741L_m$
in~\eqref{apptr:energy-recurrence} supplies the remaining site shift,
and the other final term supplies the right field with the required
positive sign.  This proves~\eqref{apptr:energy-recurrence}.

For specificity, the whole contraction can be implemented by keeping
just three rows.  Begin with
\[
 x=L_0,\qquad y=L_0T_*,\qquad w=-300000L_0G_*.
\]
For each $m=2,3,\ldots,240$, make the following updates, whose
right-hand sides all use the old rows:
\begin{align}
 w_{\mathrm{new}}&=wT_*+500000xO_*+700741(xT_*)T_*,\notag\\
 x_{\mathrm{new}}&=y,\qquad y_{\mathrm{new}}=yT_*.
 \label{apptr:algorithm}
\end{align}
After this iteration, $(x,y,w)=(L_{m-1},L_m,N_m)$.
At a required length compute
\begin{equation}\label{apptr:readout}
 V=yu,\qquad H=wu+300000x(G_*u)+700741V.
\end{equation}
Thus $(V,H)=(V_m,H_m^*)$.  The initialization realizes the same
invariant at $m=1$, which proves the indexing of the algorithm.
Ordinary exact integer matrix multiplication suffices throughout;
sparsity can reduce the work but changes no operation or comparison.

\subsection{The integer comparisons}

Table~\ref{apptr:certificate-table} gives the results of
\eqref{apptr:algorithm}--\eqref{apptr:readout}.  The integers $p_m$,
$q_m$, and $k_m$ in the table have the following exact meanings:
\begin{align}
 10^{p_m-1}&\le V_m<10^{p_m},\notag\\
 q_m10^{p_m-6}&\le V_m<(q_m+1)10^{p_m-6},\notag\\
 k_mV_m&\le20H_m^*<(k_m+1)V_m.
 \label{apptr:certificate-inequalities}
\end{align}
The last comparison is equivalent to
$k_m=\lfloor10^7H_m^*/(500000V_m)\rfloor$, since
$10^7/500000=20$.  The first comparison in particular certifies
that the vector $\psi_m$ is nonzero.

\begin{table}[htbp]
\centering
\begin{tabular}{r|r|r|r|c}
 $m$&$p_m$&$q_m$&$k_m$&
 $H_m^*/(500000V_m)$ lies in\\ \hline
 34&202&715867&4031974&$[.4031974,.4031975)$\\
 78&440&747265&4041175&$[.4041175,.4041176)$\\
 142&787&121288&4041110&$[.4041110,.4041111)$\\
 240&1317&165352&4041006&$[.4041006,.4041007)$
\end{tabular}
\caption{Exact integer contraction certificates.  The intervals
are rational intervals obtained from integer division, with strict
upper endpoints.}
\label{apptr:certificate-table}
\end{table}

All four strict comparisons
\begin{equation}\label{apptr:strict-margin}
 202100V_m-H_m^*>0
 \qquad(m=34,78,142,240)
\end{equation}
follow as well.  They are recorded directly in the full integer
output, and they also follow from the last line of
\eqref{apptr:certificate-inequalities}, since every displayed
$k_m+1$ is less than $4042000=20\cdot202100$.

The accompanying computational files make the integer comparison
fully reproducible.  The program
\path{verification/trial_verify.py} constructs the matrices from
the data in this appendix and executes the row recurrence with
arbitrary-precision integers.  The file
\path{verification/trial-exact-matrices.json} contains every
entry of $A_s$, the endpoint vectors, the restricted paired labels,
and the matrices $T_*,G_*,O_*$.  The file
\path{verification/trial-exact-results.json} contains the full
integers $V_m$, $H_m^*$, and $202100V_m-H_m^*$ for each row,
rather than only their shortened display in
Table~\ref{apptr:certificate-table}.  The equations and tables of
this appendix specify the calculation independently of these files.

\begin{proof}[Proof of Proposition~\ref{apptr:bound}]
For each of the four lengths, the integer
$V_m=\langle\psi_m,\psi_m\rangle$ is positive by
\eqref{apptr:certificate-inequalities}.  The variational principle,
Lemma~\ref{apptr:contractions}, and~\eqref{apptr:strict-margin} give
\[
 E_0^{\mathrm{opp}}(m)+am
 \le\frac{\langle\psi_m,(H_m^{\mathrm{opp}}+am\Id)\psi_m\rangle}
          {\langle\psi_m,\psi_m\rangle}
 =\frac{H_m^*}{500000V_m}
 <\frac{202100}{500000}
 =\alpha.
\]
This proves all four inequalities with the same fixed field $h=3/5$.
\end{proof}
